% SPDX-License-Identifier: MIT
% Copyright (c) 2026 Evgeny Baulin
% Optimization Methods, Seminar 10: theory handout (English).
% Generated by checks/10. Knapsack and Linear Programming in Practice.ipynb from web/10/js/content-*.js and content/students.json. Do not edit by hand.
% Compile with: latexmk -xelatex -interaction=nonstopmode -halt-on-error Theory_en.tex
\documentclass[11pt,a4paper]{article}

\usepackage{fontspec}
\setmainfont{cmunrm.otf}[BoldFont=cmunbx.otf, ItalicFont=cmunti.otf, BoldItalicFont=cmunbi.otf]
\setsansfont{cmunss.otf}[BoldFont=cmunsx.otf, ItalicFont=cmunsi.otf, BoldItalicFont=cmunso.otf]
\setmonofont{cmuntt.otf}[BoldFont=cmuntb.otf, ItalicFont=cmunit.otf, BoldItalicFont=cmuntx.otf]
\usepackage{polyglossia}
\setdefaultlanguage{english}

\usepackage[a4paper,left=2.3cm,right=2.3cm,top=2.4cm,bottom=2.4cm,headheight=14pt]{geometry}
\usepackage{amsmath,amssymb}
\usepackage{graphicx}
\usepackage{array,booktabs,longtable}
\usepackage{enumitem}
\usepackage{xcolor}
\usepackage[most]{tcolorbox}
\usepackage{titlesec}
\usepackage{fancyhdr}
\usepackage{needspace}
\usepackage[expansion=false]{microtype}
\usepackage[colorlinks=true,linkcolor=hseblue,urlcolor=hseblue,citecolor=hseblue]{hyperref}

\definecolor{hseblue}{RGB}{16,45,105}
\definecolor{hseaccent}{RGB}{198,40,40}
\definecolor{hsegray}{RGB}{108,114,126}
\definecolor{hselight}{RGB}{240,243,248}
\definecolor{hsewarm}{RGB}{253,244,232}

\hypersetup{pdftitle={Optimization Methods. Seminar 10. Theory handout}, pdfauthor={Evgeny Baulin}}

\titleformat{\section}{\color{hseblue}\Large\bfseries}{\thesection}{0.6em}{}
\titleformat{\subsection}{\color{hseblue}\large\bfseries}{\thesubsection}{0.6em}{}
\titlespacing*{\section}{0pt}{16pt plus 4pt minus 2pt}{8pt plus 2pt}
\titlespacing*{\subsection}{0pt}{12pt plus 3pt minus 2pt}{6pt plus 2pt}
\pretocmd{\section}{\Needspace{7\baselineskip}}{}{\errmessage{cannot patch section}}
\pretocmd{\subsection}{\Needspace{5\baselineskip}}{}{\errmessage{cannot patch subsection}}

\tcbset{handoutbox/.style={enhanced, breakable, sharp corners, boxrule=0pt, left=7pt, right=7pt, top=5pt, bottom=5pt, before skip=8pt plus 2pt, after skip=8pt plus 2pt}}
\newtcolorbox{examplebox}{handoutbox, colback=hselight, colframe=hselight, before upper={\textcolor{hseblue}{\bfseries Example.}\ }}
\newtcolorbox{whybox}{handoutbox, colback=white, colframe=hseblue, leftrule=2.5pt, top=3pt, bottom=3pt, before upper={\textcolor{hseaccent}{\bfseries Why.}\ }}
\newtcolorbox{answerbox}{handoutbox, colback=hsewarm, colframe=hsewarm, before upper={\textcolor{hseaccent}{\bfseries Answer.}\ }}
\newtcolorbox{conventionbox}{handoutbox, colback=hselight, colframe=hseblue, leftrule=2.5pt, before upper={\textcolor{hseblue}{\bfseries Notation.}\par\smallskip}}

\newcommand{\handoutfigure}[3]{\par\medskip\noindent\begin{minipage}{\linewidth}\centering
  \includegraphics[width=#1\linewidth]{figures/#2.pdf}\par\smallskip
  \begin{minipage}{0.94\linewidth}\small\textcolor{hseblue}{\bfseries Figure.}\ #3\end{minipage}\end{minipage}\par\medskip}
\newcommand{\problemhead}[3]{\par\bigskip\noindent{\color{hseblue}\large\bfseries #1.\ #2}\par\nopagebreak
  \noindent{\small\itshape\color{hsegray}#3}\par\nopagebreak\smallskip}
\newtcolorbox{definitionbox}[1]{handoutbox, colback=white, colframe=hseblue, leftrule=2.5pt, before upper={\textcolor{hseblue}{\bfseries #1}\ }}
\newtcolorbox{theorembox}[1]{handoutbox, colback=hselight, colframe=hseblue, leftrule=2.5pt, before upper={\textcolor{hseblue}{\bfseries #1}\ }}
\newtcolorbox{proofbox}{handoutbox, colback=white, colframe=hsegray, leftrule=1.2pt, top=3pt, bottom=3pt, before upper={\textcolor{hsegray}{\bfseries\itshape Proof.}\ }}
\newtcolorbox{notebox}[1]{handoutbox, colback=hselight, colframe=hselight, before upper={\textcolor{hseblue}{\bfseries #1}\par\smallskip}}
\newtcolorbox{trapbox}{handoutbox, colback=hsewarm, colframe=hseaccent, leftrule=2.5pt, before upper={\textcolor{hseaccent}{\bfseries Trap.}\ }}
\newtcolorbox{selfcheckbox}{handoutbox, colback=white, colframe=hseaccent, leftrule=2.5pt, top=3pt, bottom=3pt, before upper={\textcolor{hseaccent}{\bfseries Check yourself.}\ }}
\newcommand{\handoutdisplay}[1]{\sbox0{$\displaystyle #1$}\ifdim\wd0>\linewidth
  \par{\centering\binoppenalty=10000\relpenalty=0\relax$\displaystyle #1$\par}\else\[#1\]\fi}
\newcommand{\selfcheckanswer}{\par\smallskip\textcolor{hsegray}{\bfseries Answer.}\ }

\setlist[itemize]{itemsep=3pt,topsep=4pt}
\setlist[enumerate]{itemsep=3pt,topsep=4pt}
\setlength{\parindent}{0pt}
\setlength{\parskip}{5pt plus 1pt}
\setlength{\emergencystretch}{2em}
\setcounter{tocdepth}{2}

\pagestyle{fancy}
\fancyhf{}
\fancyhead[L]{\small\color{hsegray}Optimization Methods}
\fancyhead[R]{\small\color{hsegray}Seminar 10}
\fancyfoot[L]{\small\color{hsegray}Evgeny Baulin}
\fancyfoot[C]{\small\thepage}
\fancyfoot[R]{\small\color{hsegray}MIT License}
\renewcommand{\headrulewidth}{0pt}
\renewcommand{\footrulewidth}{0pt}
\fancypagestyle{plain}{\pagestyle{fancy}}

\begin{document}

\begin{center}
  {\large\color{hsegray}Optimization Methods\par}
  \vspace{4pt}
  {\large\bfseries\color{hseaccent}Seminar 10\par}
  \vspace{8pt}
  {\LARGE\bfseries\color{hseblue}The Knapsack, Its Bounds and Integer Models\par}
  \vspace{8pt}
  {\Large Theory handout\par}
  \vspace{12pt}
  HSE University, Faculty of Computer Science\par
  Academic year 2026-2027\par
  \vspace{8pt}
  Course lecturer: Andrey Ignatov\par
  Evgeny Baulin\par
  Telegram: \href{https://t.me/tarakan_tuc}{@tarakan\_tuc}\par
\end{center}

\noindent{\color{hseblue}\rule{\linewidth}{0.6pt}}

\tableofcontents

\clearpage


This handout follows seminar 10. It starts from the 0–1 knapsack and its relaxation, builds branch and bound and dynamic programming on it, then turns to modelling with binary variables. Every problem of the seminar is here with its full solution.

\begin{conventionbox}
Maximization unless stated otherwise. Profits $p_i>0$, weights $w_i>0$, capacity $W$; $x_i\in\{0,1\}$ says whether item $i$ is packed. $U(N)$ is an upper bound for a node $N$ of the search, $z_{\text{inc}}$ the value of the best feasible point found so far. $D(i,c)$ is the best profit from the first $i$ items in capacity $c$.
\end{conventionbox}

% =====================================================================
\section{The knapsack and the density bound}\label{sec:knapsack}

The knapsack is the simplest integer program with one constraint, and almost everything about integer optimization can be seen on it: bounds, search, tables and the gap between the relaxation and the integers.

\begin{definitionbox}{Definition: 0–1 knapsack and its relaxation.}
$\max\sum_ip_ix_i$ subject to $\sum_iw_ix_i\le W$, $x\in\{0,1\}^n$. The relaxation keeps the objective and the row and replaces $x_i\in\{0,1\}$ by $0\le x_i\le1$. Variants: bounded copies $0\le x_i\le u_i$ integer, unlimited copies $x_i\in\mathbb Z_{\ge0}$, several resource rows, several knapsacks.
\end{definitionbox}

\begin{theorembox}{Theorem: Fractional knapsack.}
Sort the items by density $p_i/w_i$, decreasing. Take whole items while they fit, then the fraction of the next item that fills the capacity. This solves the relaxation, and at most one variable is fractional.
\end{theorembox}

\begin{proofbox}
Take an optimal relaxed solution in which a denser item $i$ is not full while a less dense item $j$ uses weight. Moving a small weight $\delta$ from $j$ to $i$ keeps the row and changes the profit by $\delta(p_i/w_i-p_j/w_j)\ge0$. Repeating this ends with a full prefix of the sorted list and one partial item, which is the greedy solution.
\end{proofbox}

\begin{whybox}
The relaxation is a sale by weight: you pay with capacity, so you buy the goods with the best price per kilogram first. Integrality is what breaks this: the last item cannot be cut.
\end{whybox}

\handoutfigure{0.72}{fig_a1}{The drone of A1: its capacity taken item by item in density order, and the optimal load underneath.}

\problemhead{A1}{Parcels for a drone}{On the board}
\textbf{Statement.} Five parcels pay $(17,\ 12,\ 7,\ 19,\ 17)$ for delivery and weigh $(8,\ 6,\ 7,\ 9,\ 4)$ kg; the drone lifts at most $W=19$ kg. Solve the relaxation by densities, round the fractional parcel down, and find the best load. Give the bound $U$, the value of the rounded load and the optimum.

\textbf{Solution.}
\begingroup\sloppy
\begin{enumerate}[label=\textbf{\arabic*.}, leftmargin=*]
  \item The densities are $(\tfrac{17}{8},\ 2,\ 1,\ \tfrac{19}{9},\ \tfrac{17}{4})$, so the order is $(5,\ 1,\ 4,\ 2,\ 3)$. Parcels are taken whole in this order until one no longer fits.
  \item Parcels 5 and 1 fit; parcel 4 does not, and the remaining capacity takes the fraction $\tfrac{7}{9}$ of it. The relaxed solution is $x=(1,0,0,\tfrac{7}{9},1)$ with $U=\tfrac{439}{9}\approx48.7778$.
  \item Rounding $x_{ 4 }$ down leaves a load worth $34$ and a lot of free capacity. Skipping the parcel that does not fit and going on down the list gives $46$ with parcels $(5,\ 1,\ 2)$: better, still not optimal.
  \item Parcels $(2,\ 4,\ 5)$ weigh exactly $W$ and pay $48$. The profits are integers, so no load pays more than $\lfloor U\rfloor$, which is $48$ as well. The bound proves that this load is optimal.
\end{enumerate}
\endgroup
\begin{answerbox}
$U=\tfrac{439}{9}\approx48.7778$, the rounded load $34$, the optimum $48$ with parcels $(2,\ 4,\ 5)$.
\end{answerbox}

\begin{trapbox}
Rounding the relaxed solution down is a heuristic, not a method: in A1 it loses almost a third of the value. With integer profits, $\lfloor U\rfloor$ is the honest upper bound.
\end{trapbox}

\problemhead{A2}{Weight and volume}{On your own}
\textbf{Statement.} Five boxes pay $(13,\ 7,\ 11,\ 16,\ 9)$, weigh $(7,\ 6,\ 5,\ 9,\ 3)$ and take volumes $(6,\ 2,\ 6,\ 6,\ 3)$; the van carries weight $W=17$ and volume $V=10$. Take boxes in the order of profit per unit of weight while both limits allow, then find the best load and the value of the relaxation with both rows. Give the greedy value, the optimum and the LP bound.

\textbf{Solution.}
\begingroup\sloppy
\begin{enumerate}[label=\textbf{\arabic*.}, leftmargin=*]
  \item By profit per unit of weight the order is $(5,\ 3,\ 1,\ 4,\ 2)$. Going down it and skipping what breaks a limit, the greedy load is $(5,\ 3)$ worth $20$: the light dense boxes use up the volume.
  \item Checking the packings that respect both limits, the best is $(4,\ 5)$ worth $25$.
  \item The relaxation with both rows has its optimum at $x=(0,1,0,\tfrac{5}{6},1)$ with value $\tfrac{88}{3}\approx29.3333$. The volume row is tight there and the weight row still has room, so only one variable is fractional; with both rows tight a vertex could have two.
\end{enumerate}
\endgroup
\begin{answerbox}
Greedy $20$, optimum $25$ with boxes $(4,\ 5)$, LP bound $\tfrac{88}{3}\approx29.3333$.
\end{answerbox}

\problemhead{A3}{Copies of the same item}{On your own}
\textbf{Statement.} Three kinds of item pay $(7,\ 14,\ 12)$ and weigh $(7,\ 3,\ 2)$, with at most $u=4$ copies of each kind; the capacity is $W=13$. Solve the relaxation $0\le x_i\le u$, then find the integer optimum. Give the LP bound, the optimum and the number of copies of the second kind in it.

\textbf{Solution.}
\begingroup\sloppy
\begin{enumerate}[label=\textbf{\arabic*.}, leftmargin=*]
  \item The third kind is the densest, then the second. The relaxation takes $u$ copies of the third and the fraction of the second that fills the rest: $U=\tfrac{214}{3}\approx71.3333$.
  \item Rounding that down and filling greedily in whole copies gives $62$. Two copies fewer of the densest kind leave room for two more of the second: the optimum is $x=(0,\ 3,\ 2)$ with value $66$, and the capacity is used to the last unit.
\end{enumerate}
\endgroup
\begin{answerbox}
$U=\tfrac{214}{3}\approx71.3333$, optimum $66$ at $x=(0,\ 3,\ 2)$.
\end{answerbox}

\begin{selfcheckbox}
Why does sorting by density fail to solve the relaxation with two resource rows?
\selfcheckanswer There is no single price per unit: an item can be light and bulky. The relaxation is a real LP, and its vertex can have two fractional variables.
\end{selfcheckbox}

% =====================================================================
\section{Branch and bound}\label{sec:bnb}

Enumeration of all $2^n$ packings is a binary tree: each level fixes one variable. Branch and bound walks this tree but skips every subtree that provably holds nothing better than what it already has.

\begin{definitionbox}{Definition: Node, bound, incumbent.}
A node $N$ fixes some variables; its completions are the packings that agree with them. $U(N)$ is a valid upper bound if every feasible completion has value at most $U(N)$; for the knapsack it is the profit of the items fixed at one plus the relaxed value of the other free items in the capacity that is left. The incumbent $z_{\text{inc}}$ is the best feasible value found so far.
\end{definitionbox}

\begin{theorembox}{Theorem: Pruning and exactness.}
If $U(N)\le z_{\text{inc}}$, the node $N$ contains no better solution. With valid bounds and complete branching, the incumbent is optimal when no node is left open.
\end{theorembox}

\begin{proofbox}
Any feasible completion $x$ of $N$ has $f(x)\le U(N)\le z_{\text{inc}}$, so discarding $N$ loses nothing. Every branch fixes one more variable, so the tree is finite, and every feasible point lies in exactly one leaf that was either examined or discarded for this reason.
\end{proofbox}

\begin{itemize}
  \item A node is finished in three cases: its fixed items exceed the capacity; $U(N)\le z_{\text{inc}}$; the relaxed optimum at the node is already a packing, which may replace the incumbent.
  \item Branch on the fractional item of the relaxation: both children change the bound, and one of them often closes at once.
  \item Depth first: little memory, an incumbent soon. Best bound: fewer nodes to prove optimality, more open nodes in memory.
  \item An unfinished run still certifies something: the incumbent bounds the optimum from below, the biggest bound among the open nodes from above.
\end{itemize}

\handoutfigure{0.66}{fig_b1}{The search tree of B1 with every bound: the cross is the infeasible node, the orange one is closed by its bound, the red ones are incumbents.}

\problemhead{B1}{A search tree by hand}{On the board}
\textbf{Statement.} Profits $(9,\ 14,\ 15,\ 14)$, weights $(6,\ 3,\ 4,\ 9)$, capacity $W=13$. Run branch and bound depth first. At every node compute the density bound, branch on the fractional item with the branch $x_j=1$ first, and close a node that is infeasible, integral, or has $U(N)\le z_{\text{inc}}$. Count every node whose bound you compute, the root and the infeasible ones included. Give the number of nodes and the optimum.

\textbf{Solution.}
\begingroup\sloppy
\begin{enumerate}[label=\textbf{\arabic*.}, leftmargin=*]
  \item Root: items 2 and 3 fit, item 4 only in part, so $U=\tfrac{115}{3}\approx38.3333$ and we branch on $x_{ 4 }$.
  \item With $x_{ 4 }=1$ the bound is $31.75$ and item 3 is fractional. The branch $x_{ 3 }=1$ fills the knapsack exactly with an integral relaxation: the first incumbent, $z_{\text{inc}}=29$.
  \item The branch $x_{ 3 }=0$ has $U=29.5$, above the incumbent, so we branch on item 1. Taking it overloads the knapsack (infeasible); leaving it gives $U=28\le z_{\text{inc}}$, closed by the bound.
  \item Back at the root, $x_{ 4 }=0$: the other three items fill the knapsack exactly, the relaxation is integral with value $38$, a new incumbent. Nothing is open, so it is optimal. The whole log:
\handoutdisplay{
\begin{array}{c|l|c|c} k & \text{fixed} & U & z_{\text{inc}} \\ \hline 1 & \text{---} & 38.33 & 0 \\ 2 & x_4=1 & 31.75 & 0 \\ 3 & x_4=1,\ x_3=1 & 29 & 0 \\ 4 & x_4=1,\ x_3=0 & 29.5 & 29 \\ 5 & x_4=1,\ x_3=0,\ x_1=1 & \text{---} & 29 \\ 6 & x_4=1,\ x_3=0,\ x_1=0 & 28 & 29 \\ 7 & x_4=0 & 38 & 29 \end{array}
}

\end{enumerate}
\endgroup
\begin{answerbox}
7 nodes, the optimum $38$ with items $(1,\ 2,\ 3)$.
\end{answerbox}

\problemhead{B2}{Depth first or best bound}{On your own}
\textbf{Statement.} Profits $(6,\ 7,\ 17,\ 19,\ 20)$, weights $(7,\ 4,\ 6,\ 5,\ 9)$, capacity $W=18$. Run branch and bound with the rules of B1 twice: depth first, and best bound, which always expands the open node with the largest $U$; there a child is closed as soon as its bound is computed if it is infeasible, integral or not above the incumbent. Count the nodes whose bound you compute, the root included. Give both counts and the optimum.

\textbf{Solution.}
\begingroup\sloppy
\begin{enumerate}[label=\textbf{\arabic*.}, leftmargin=*]
  \item The density order is $(4,\ 3,\ 5,\ 2,\ 1)$ and the root bound is $\tfrac{464}{9}$. Depth first goes down the branches $x_j=1$ and meets three infeasible children before the first packing, worth the first incumbent in the log:
\handoutdisplay{
\begin{array}{c|l|c|c} k & \text{fixed} & U & z_{\text{inc}} \\ \hline 1 & \text{---} & 51.56 & 0 \\ 2 & x_5=1 & 50.33 & 0 \\ 3 & x_5=1,\ x_3=1 & 48.4 & 0 \\ 4 & x_5=1,\ x_3=1,\ x_4=1 & \text{---} & 0 \\ 5 & x_5=1,\ x_3=1,\ x_4=0 & 42.25 & 0 \\ 6 & x_5=1,\ x_3=1,\ x_4=0,\ x_2=1 & \text{---} & 0 \\ 7 & x_5=1,\ x_3=1,\ x_4=0,\ x_2=0 & 39.57 & 0 \\ 8 & x_5=1,\ x_3=1,\ x_4=0,\ x_2=0,\ x_1=1 & \text{---} & 0 \\ 9 & x_5=1,\ x_3=1,\ x_4=0,\ x_2=0,\ x_1=0 & 37 & 0 \\ 10 & x_5=1,\ x_3=0 & 46 & 37 \\ 11 & x_5=0 & 45.57 & 46 \end{array}
}
 That is 11 nodes.
  \item Best bound expands the root's child with the larger bound, and one of its children is already a packing worth $46$. It then expands the only open node still above that value; its children are infeasible or below it, and the root's other child, below the incumbent, is never expanded. It computes 7 bounds.
  \item Both end with the optimum $46$, items $(2,\ 4,\ 5)$.
\end{enumerate}
\endgroup
\begin{answerbox}
Depth first 11 nodes, best bound 7, optimum $46$.
\end{answerbox}

\begin{trapbox}
Counting nodes needs a convention: whether the root, the infeasible children and the children closed at once are counted. Compare strategies only under the same convention.
\end{trapbox}

% =====================================================================
\section{Cutting planes}\label{sec:cuts}

Instead of splitting the problem, one can tighten it: add an inequality that holds at each integer point and fails at the relaxed optimum found so far. Solvers do both, which is called branch and cut.

\begin{theorembox}{Theorem: Gomory's fractional cut.}
Let the data be integers and let a row of the optimal tableau read $x_k+\sum_{j\in N}\bar a_jx_j=\bar b$ with $\bar b$ fractional, $N$ the nonbasic variables. Then $\sum_{j\in N}\{\bar a_j\}x_j\ge\{\bar b\}$ holds at every integer point and fails at the LP optimum, where $\{t\}=t-\lfloor t\rfloor$.
\end{theorembox}

\begin{proofbox}
Split each coefficient into its integer and fractional parts: $x_k+\sum\lfloor\bar a_j\rfloor x_j-\lfloor\bar b\rfloor=\{\bar b\}-\sum\{\bar a_j\}x_j$. At an integer point the left side is an integer, and the right side is below one, so it is at most zero. At the LP optimum all $x_j$, $j\in N$, are zero and the right side is $\{\bar b\}>0$.
\end{proofbox}

\begin{whybox}
The cut slices off the corner of the polytope that holds the fractional optimum and nothing else that matters. In two variables one can see it: the line $x_1+x_2\le d$ of B3 passes through integer points.
\end{whybox}

\problemhead{B3}{One cut of Gomory}{On your own}
\textbf{Statement.} Maximize $4x_1+2x_2$ subject to $6x_1+2x_2\le14$, $x_1+2x_2\le10$, $x\ge0$ integer. Solve the relaxation, take the row of $x_1$ in the optimal tableau and derive Gomory's cut; written in $x$, it reads $x_1+x_2\le d$. Give the LP value, $d$ and the LP value after adding the cut.

\textbf{Solution.}
\begingroup\sloppy
\begin{enumerate}[label=\textbf{\arabic*.}, leftmargin=*]
  \item Both rows tight give $x=(\tfrac{4}{5},\tfrac{23}{5})$ with $z=\tfrac{62}{5}=12.4$; both variables are fractional.
  \item With $x_1,x_2$ basic the row of $x_1$ is
\handoutdisplay{
x_1+\tfrac{1}{5}s_1-\tfrac{1}{5}s_2=\tfrac{4}{5}
}
 The fractional parts of $-\tfrac{1}{5}$ and $\tfrac{1}{5}$ are $\tfrac{4}{5}$ and $\tfrac{1}{5}$, so the cut is
\handoutdisplay{
\tfrac{1}{5}s_1+\tfrac{4}{5}s_2\ge\tfrac{4}{5}
}

  \item Substitute $s_1=b_1-6x_1-2x_2$ and $s_2=b_2-x_1-2x_2$: the $x$-terms add up to $-2x_1-2x_2$ and the inequality becomes $x_1+x_2\le5$. The current optimum has $x_1+x_2>5$, so it is cut off, while every integer point satisfies the cut.
  \item With the cut the LP optimum moves to $(1,\ 4)$ with value $12$. That point is integral, so it is the integer optimum: one cut closed the gap.
\end{enumerate}
\endgroup
\begin{answerbox}
$z=\tfrac{62}{5}=12.4$, the cut $x_1+x_2\le5$, after it $z=12$ at $(1,\ 4)$.
\end{answerbox}

% =====================================================================
\section{Dynamic programming for the knapsack}\label{sec:dp}

With integer weights the knapsack has a second exact method. It does not search a tree; it fills a table of subproblems, each solved once from two smaller ones.

\begin{definitionbox}{Definition: State.}
$D(i,c)$ is the best profit using only the first $i$ items in integer capacity $c$, for $0\le i\le n$ and $0\le c\le W$. The boundary values are $D(0,c)=0$, and the answer is $D(n,W)$.
\end{definitionbox}

\begin{theorembox}{Theorem: The recurrence.}
$D(i,c)=D(i-1,c)$ if $w_i>c$, and $D(i,c)=\max\{D(i-1,c),\ p_i+D(i-1,c-w_i)\}$ otherwise.
\end{theorembox}

\begin{proofbox}
An optimal choice among the first $i$ items either leaves item $i$ out, and is then optimal for $(i-1,c)$, or takes it, and then the rest is optimal for $(i-1,c-w_i)$. Both options are achievable, so the larger of the two is optimal. Induction on $i$ finishes the proof.
\end{proofbox}

\begin{itemize}
  \item The way back: start at $(n,W)$; when $D(i,c)=D(i-1,c)$ skip item $i$, else take it and subtract $w_i$ from $c$.
  \item One row of memory suffices for the value: update it from $c=W$ down to $c=w_i$, so that each item is used once.
  \item Unlimited copies: $D(i,c)=\max\{D(i-1,c),\ p_i+D(i,c-w_i)\}$, the same row; in one row, scan upwards.
  \item Bounded copies: maximize over $q=0,\dots,\min(u_i,\lfloor c/w_i\rfloor)$ copies of item $i$.
\end{itemize}

\handoutfigure{0.72}{fig_c1}{C1 as a table; the framed cells trace the packing back.}

\problemhead{C1}{The table of a small knapsack}{On the board}
\textbf{Statement.} Profits $(13,\ 10,\ 6,\ 5)$, weights $(5,\ 4,\ 3,\ 2)$, capacity $W=8$. Fill the table $D(i,c)$ for $i=0,\dots,4$ and $c=0,\dots,W$, read the optimum and walk back to the packing. Give $D(2,W)$ and $D(4,W)$.

\textbf{Solution.}
\begingroup\sloppy
\begin{enumerate}[label=\textbf{\arabic*.}, leftmargin=*]
  \item Row $0$ is all zeros. Row $1$ jumps to $p_1$ at $c=w_1$; row $2$ adds item $2$, but items $1$ and $2$ together are too heavy, so $D(2,W)=13$.
  \item Rows $3$ and $4$ by the same rule give the whole table:
\handoutdisplay{
\begin{array}{c|ccccccccc} i\backslash c & 0 & 1 & 2 & 3 & 4 & 5 & 6 & 7 & 8 \\ \hline 0 & 0 & 0 & 0 & 0 & 0 & 0 & 0 & 0 & 0 \\ 1 & 0 & 0 & 0 & 0 & 0 & 13 & 13 & 13 & 13 \\ 2 & 0 & 0 & 0 & 0 & 10 & 13 & 13 & 13 & 13 \\ 3 & 0 & 0 & 0 & 6 & 10 & 13 & 13 & 16 & 19 \\ 4 & 0 & 0 & 5 & 6 & 10 & 13 & 15 & 18 & 19 \end{array}
}
 The corner is $D(4,W)=19$.
  \item Way back from $(4,W)$: $D(4,W)=D(3,W)$, so item $4$ is out; $D(3,W)\ne D(2,W)$, so item $3$ is in and $c$ drops by $w_3$; item $2$ is out and item $1$ is in. The packing is $(1,\ 3)$.
  \item The density order would take item $1$, skip item $2$ and add item $4$: $18$, one less than the optimum. The table has $(n+1)(W+1)=45$ cells, each filled in constant time.
\end{enumerate}
\endgroup
\begin{answerbox}
$D(2,W)=13$, $D(4,W)=19$ with items $(1,\ 3)$.
\end{answerbox}

\problemhead{C2}{Unlimited copies}{On your own}
\textbf{Statement.} Three kinds of item, as many copies as you like, pay $(14,\ 17,\ 11)$ and weigh $(4,\ 5,\ 3)$; the capacity is $W=13$. Fill one row $D(c)$ for $c=0,\dots,W$ with $D(c)=\max\bigl(D(c-1),\ \max_{i:\,w_i\le c}\{p_i+D(c-w_i)\}\bigr)$, then compare with taking the densest kind as often as it fits and filling the rest greedily. Give $D(W)$, the number of copies of the third kind in the optimum and the greedy value.

\textbf{Solution.}
\begingroup\sloppy
\begin{enumerate}[label=\textbf{\arabic*.}, leftmargin=*]
  \item The row is $D=(0,\ 0,\ 0,\ 11,\ 14,\ 17,\ 22,\ 25,\ 28,\ 33,\ 36,\ 39,\ 44,\ 47)$ for $c=0,1,\dots,W$. The corner is $D(W)=47$.
  \item Walking back by the winning items gives the copies $x=(1,\ 0,\ 3)$: they use the whole capacity.
  \item The third kind is the densest. Taking it as often as it fits leaves a remainder that no item fills, and the greedy load is worth $44$.
\end{enumerate}
\endgroup
\begin{answerbox}
$D(W)=47$ with $x=(1,\ 0,\ 3)$; greedy $44$.
\end{answerbox}

\problemhead{C3}{Exact totals}{On your own}
\textbf{Statement.} Weights $(5,\ 6,\ 9,\ 14)$, each used at most once. Mark which totals $c=0,\dots,W$ with $W=21$ can be made exactly, with one row of booleans updated for each weight from $c=W$ down to $c=w_i$. Give the number of reachable totals, $0$ included, and the largest one. Is the optimum of the knapsack with $p_i=w_i$ equal to $W$?

\textbf{Solution.}
\begingroup\sloppy
\begin{enumerate}[label=\textbf{\arabic*.}, leftmargin=*]
  \item After all four weights the reachable totals up to $W$ are $(0,\ 5,\ 6,\ 9,\ 11,\ 14,\ 15,\ 19,\ 20)$: $9$ of them.
  \item The largest is $20$, short of $W$, so no subset sums to $W$ and the knapsack with $p_i=w_i$ has optimum $20$, not $W$.
  \item The row has $W+1$ booleans and each weight scans it once: $O(nW)$ again, which is why subset sum, an NP-complete problem, is easy for small $W$.
\end{enumerate}
\endgroup
\begin{answerbox}
$9$ reachable totals, the largest $20$; the optimum is not $W$.
\end{answerbox}

\begin{trapbox}
$O(nW)$ is not polynomial: $W$ is written with about $\log_2W$ bits, so the table is exponential in the input length. Multiplying all weights and the capacity by a thousand leaves the problem the same and makes the table about a thousand times larger.
\end{trapbox}

\begin{selfcheckbox}
When would you choose branch and bound over the table?
\selfcheckanswer When the weights are large or not integers, when the bounds are strong, or when only one capacity matters. The table wins with small integer capacities or when every capacity is needed.
\end{selfcheckbox}

% =====================================================================
\section{Binary variables in models}\label{sec:modelling}

Most practical integer models are linear programs with a few binary switches. The switches carry logic, fixed costs and interactions; the art is to write them so that the relaxation stays close to the integers.

\begin{itemize}
  \item If $A$ then $B$: $y_A\le y_B$. Not both: $y_A+y_B\le1$. At most $k$ of a set: $\sum y_i\le k$. Exactly one: $\sum y_i=1$.
  \item Fixed charge: $x\le u\,y$ with a valid bound $u$, and $-f\,y$ in the objective. The smaller the valid $u$, the stronger the relaxation.
  \item Product of binaries $z=x_ix_j$: $z\le x_i$, $z\le x_j$, $z\ge x_i+x_j-1$, $z\ge0$; keep only the side the objective pushes against.
  \item Binary times continuous $z=y\,x$ with $0\le x\le u$: $z\le u\,y$, $z\le x$, $z\ge x-u(1-y)$, $z\ge0$.
\end{itemize}

\problemhead{D1}{Three lines and two rules}{On the board}
\textbf{Statement.} Three lines make products with profit per unit $(7,\ 8,\ 6)$; opening a line costs $(40,\ 20,\ 40)$. A unit takes $(3,\ 4,\ 1)$ hours of a shared crew that has $80$ hours, and demand caps the output at $(15,\ 10,\ 25)$. Line 1 may open only if line 2 is open, and at most two lines may open. Write the model with $y_j\in\{0,1\}$ and $x_j\ge0$, and solve it by trying the allowed $y$. Give the best profit without the two rules and with them.

\textbf{Solution.}
\begingroup\sloppy
\begin{enumerate}[label=\textbf{\arabic*.}, leftmargin=*]
  \item The model: $\max\sum_jp_jx_j-\sum_jf_jy_j$ subject to $\sum_ja_jx_j\le H$, $x_j\le u_jy_j$, $y_1\le y_2$, $y_1+y_2+y_3\le2$, $x\ge0$, $y\in\{0,1\}^3$. The row $x_j\le u_jy_j$ links the output to the line and the demand at once.
  \item Without the rules, lines $1$ and $3$ are best: output $(15,\ 0,\ 25)$ and profit $175$. This plan breaks the first rule.
  \item The patterns that obey both rules, each with its best output:
\handoutdisplay{
\begin{array}{ccc|ccc|c} y_1 & y_2 & y_3 & x_1 & x_2 & x_3 & \text{profit} \\ \hline 0 & 0 & 0 & 0 & 0 & 0 & 0 \\ 0 & 0 & 1 & 0 & 0 & 25 & 110 \\ 0 & 1 & 0 & 0 & 10 & 0 & 60 \\ 0 & 1 & 1 & 0 & 10 & 25 & 170 \\ 1 & 1 & 0 & 15 & 8.75 & 0 & 115 \end{array}
}
 The best is lines $2$ and $3$, $x=(0,\ 10,\ 25)$, profit $170$.
\end{enumerate}
\endgroup
\begin{answerbox}
Without the rules $175$, with them $170$ (lines $2$ and $3$).
\end{answerbox}

\problemhead{D2}{Neighbouring stands}{On your own}
\textbf{Statement.} Four stands have values $(10,\ 4,\ 5,\ 5)$ and costs $(2,\ 2,\ 2,\ 5)$, with budget $B=11$. Neighbours interact: the pairs $(1,2)$, $(2,3)$ and $(3,4)$ add $(3,\ -8,\ 3)$ to the value when both stands of the pair are chosen. Replace each product $x_ix_{i+1}$ by a new variable with only the rows it needs, and solve. Give the optimum, the value the model reports if the pair $(2,3)$ gets only $z\le x_2$ and $z\le x_3$, and the number of linearizing rows you used.

\textbf{Solution.}
\begingroup\sloppy
\begin{enumerate}[label=\textbf{\arabic*.}, leftmargin=*]
  \item The two bonuses get $z\le x_i$ and $z\le x_{i+1}$; the objective pushes them up, so they equal $1$ exactly when both stands are chosen. The penalty gets $z_{2,3}\ge x_2+x_3-1$; the objective pushes it down to $0$ unless both are chosen. That is $5$ rows.
  \item Checking the subsets within the budget, the best is $(1,\ 3,\ 4)$ with value $23$: stand $2$ is dropped to avoid the penalty.
  \item With only $z_{2,3}\le x_2$, $z_{2,3}\le x_3$ the solver is free to set $z_{2,3}=0$ even when both stands are chosen, so the penalty is never paid. It picks $(1,\ 2,\ 3,\ 4)$ and reports $30$, a value no real plan reaches.
\end{enumerate}
\endgroup
\begin{answerbox}
Optimum $23$ with stands $(1,\ 3,\ 4)$; with the wrong rows $30$; $5$ rows.
\end{answerbox}

\begin{trapbox}
A missing half of a linearization raises no error. The solver simply exploits the gap and reports a value that no real plan reaches, as in D2.
\end{trapbox}

% =====================================================================
\section{Strong and weak formulations}\label{sec:formulations}

\begin{definitionbox}{Definition: Stronger formulation.}
Two formulations of the same set of integer points are compared by their relaxations: $P_1$ is stronger than $P_2$ if $P_1\subseteq P_2$. The ideal formulation is the convex hull of the integer points, where every vertex is integral.
\end{definitionbox}

\begin{theorembox}{Proposition: Bounds of nested relaxations.}
If $P_1\subseteq P_2$, then for a maximization $\max_{P_1}c^{\top}x\le\max_{P_2}c^{\top}x$: the stronger formulation gives a bound at least as good, and branch and bound closes more nodes with it.
\end{theorembox}

\problemhead{D3}{Two models of one warehouse problem}{On your own}
\textbf{Statement.} Two warehouses with opening costs $(6,\ 8)$ serve four shops; serving shop $i$ from warehouse $j$ costs $c_{ij}$, with $C=\begin{pmatrix} 7 & 2 \\ 6 & 5 \\ 2 & 8 \\ 2 & 4 \end{pmatrix}$. Each shop is served by one open warehouse. Write the model with the aggregated rows $\sum_ix_{ij}\le4y_j$ and with the rows $x_{ij}\le y_j$, solve both relaxations and the integer problem. Give the two LP values and the optimum.

\textbf{Solution.}
\begingroup\sloppy
\begin{enumerate}[label=\textbf{\arabic*.}, leftmargin=*]
  \item Aggregated: every shop goes to its cheaper warehouse, and $y=(0.5,\ 0.5)$ just covers them, paying a fraction of each opening cost. The LP value is $18$.
  \item Disaggregated: a shop served half from a warehouse forces half of that warehouse open, and fractional $y$ no longer saves anything. The LP value is $23$, and its optimum is already integral.
  \item The integer optimum opens warehouse $(1)$ only and costs $23$, equal to the stronger bound.
\end{enumerate}
\endgroup
\begin{answerbox}
LP values $18$ and $23$, optimum $23$.
\end{answerbox}

\begin{whybox}
The aggregated row $\sum_ix_{ij}\le m\,y_j$ lets a warehouse serve some shops fully while paying only a fraction of its opening cost; the rows $x_{ij}\le y_j$ say what integrality says anyway, that a shop can only use an open warehouse.
\end{whybox}

% =====================================================================
\section{The cat sudoku}\label{sec:cats}

The puzzle of block G is a feasibility problem in binaries: one cat per row, column and region, no two touching. Its solver on the page is a depth-first search over rows that closes a node as soon as a row or a region has no cell left, that is, branch and bound with infeasibility as the only bound.

Written as an integer program, it has the rows of the model in block G. The square inequalities are the stronger choice: a solver with them usually needs no branching at all on small boards.

\end{document}
